\documentclass[a4paper]{article}
\usepackage[boxedEnv, titleFormat]{myBelovedPackage}

\margin{3cm}{4cm}


% Vraisemblance
\newcommand{\Lmarg}[1][]{\mc L_{marg#1}}
\newcommand{\logmarg}[1][]{\ell_{marg#1}}

\newcommand{\Lcomp}[1][]{\mc L_{comp#1}}
\newcommand{\logcomp}[1][]{\ell_{comp#1}}

\newcommand{\varobs}{\mc O}
\newcommand{\varlat}[1][]{Z#1}
\newcommand{\pen}{\mathrm{pen}}
\DeclareMathOperator*{\prox}{Prox}


\newcommand{\hazard }{h}
\newcommand{\regularization }{\lambda}
\newcommand{\nonlinearfct}{m}


\newcommand{\forany}[3][1]{#1\leq #2 \leq #3}
%Indice
\newcommand{\geno}{g}
\newcommand{\maxgeno}{G}
\newcommand{\foranygeno}{\forany{\geno}{\maxgeno}}

\newcommand{\lines}{\ell}
\newcommand{\maxlines}{L}
\newcommand{\foranylines}{\forany{\lines}{\maxlines}}

\newcommand{\genolines}{ {\geno,\lines} }
\newcommand{\meltgenolines}{ {\geno\lines} }

\newcommand{\obs}{j}
\newcommand{\maxobs}{J}
\newcommand{\foranyobs}{\forany{\obs}{\maxobs}}

\newcommand{\genolinesobs}{ {\geno, \lines, \obs} }
\newcommand{\genoobs}{ {\geno, \obs} }




% --- lazy math command --- %
\newcommand{\psinorm}[2]{\frac 12 \log(2\pi#1^2) + \frac{#2^2}{2#1^2}}
\newcommand{\psibar}[1]{\psinorm{\bar #1}{\sigma_{#1}}}



% Petite commande pour faire des model exp = h + <S, Phi> + psi parce que je suis feignant
\newcounter{Sc}
\setcounter{Sc}{0}
\newcommand{\Scount}{{\stepcounter{Sc}\theSc}}

\newcounter{Phic}
\setcounter{Phic}{0}
\newcommand{\Phicount}{{\stepcounter{Phic}\thePhic}}

\newcounter{hc}
\setcounter{hc}{0}
\newcommand{\hcount}{{\stepcounter{hc}\thehc}}

\newcounter{psic}
\setcounter{psic}{0}
\newcommand{\psicount}{{\stepcounter{psic}\thepsic}}
    

% === Couleur === %
\setmyColor[ blue = {0, 0, 255},
             violet = {163, 73, 164},
             fg = {0, 163, 166}]

\xdefinecolor{inraeRed}{RGB}{237, 110, 108}
\xdefinecolor{inraeBlue}{RGB}{66, 48, 137}
\xdefinecolor{inraeCyan}{RGB}{0, 163, 166}
\xdefinecolor{inraeGreen}{RGB}{157, 197, 68}

\newcommand{\inrae}[1]{\tc{inraeCyan}{\bf{#1}}}
\newcommand{\known}[1]{ \tc{BrickRed}{#1}}

\newcommand{\by}[1]{\textbf{[#1]}}


% === enumerate === %
\newenvironment{myEnumerate}[2][white]{
    \setlist[enumerate]{%wide=0pt,
                        %labelsep=0.5em,
                        font=\large\color{#1}}
    
    \setlist[enumerate,1]{label = \colorbox{#2}{\makebox[0.75em][c]{\arabic*}}}
    \setlist[enumerate,2]{label = \colorbox{#2}{\makebox[0.75em][c]{\alph*}}}
    \setlist[enumerate,3]{label = \colorbox{#2}{\makebox[0.75em][c]{\roman*}}}

    \begin{enumerate}
    }{ \end{enumerate} }

\makeatletter
\newif\if@ShareVersion 

\@ShareVersiontrue
\@ShareVersionfalse

\newcommand{\remark}[1]{\if@ShareVersion \else \tc{inraeRed}{#1} \fi}


\newcommand{\myRemarks}[1]{
    \if@ShareVersion        
    \else 
        { 
            \color{inraeRed}
            \begin{myEnumerate}{inraeRed}
                #1
            \end{myEnumerate}
        }
    \fi}
    
\makeatother

\begin{document} 
\begin{myText}\centering

{\Large \underline{Écrire des mails}} \\
\includegraphics[width = 0.6\textwidth]{logos/pour_les_nuls.png}


\vspace{2cm}

\begin{myEnumerate}{inraeCyan}
    \item Pas envoyer de mails en fin de journée, \underline{traiter les mails dans la matiné !}
    \item D'abord écrire, puis laisser un peu de temps avant de l'envoyer : 
    \\pour éviter d'envoyer un mails dans le feu de l'action
    \item relire le message en se mettant à la place du destinataire pour être sûr que le message est bien clair 
    \item Solliciter Estelle et/ou Sarah en cas de doutes/mails important,
    \item Ne pas mettre le monde entier en copie,
    \item Vérifier que je m'adresse à la bonne personnes,
    \\ Vérifier que je n'ai pas sauter un échelon dans la hierachie
    \item \underline{Structure du mails :}
    \begin{myEnumerate}{inraeBlue}
        \item Bonjour prénom si je connais la personne, monsieur madame sinon

        \item remercier si j'ai demandé quelque chose
        \item Etre claire et concis autant faire ce peut.

        \item Terminer par un formulation  : cordialement/en vous remerciant, bonne journée etc etc en fonction de la hierachie ou si je connais la personne
    \end{myEnumerate}

\end{myEnumerate}

\end{myText}
\end{document}